%% ============================================================
%%  Aula 3 — Avaliação de Política Sem Modelo
%%  Curso de Aprendizagem por Reforço (pt-BR)
%%  Adaptado e expandido a partir de: Emma Brunskill, CS234
%%  (Stanford), Lecture 3. Estrutura inspirada também em
%%  David Silver (Lecture 4) e Sutton & Barto (cap. 5–6).
%%  Compilar com XeLaTeX.
%% ============================================================
\documentclass[aspectratio=169, 11pt]{beamer}
%% .sty do projeto na raiz do repositório (../)
\makeatletter
\def\input@path{{./}{../}}
\makeatother


\usetheme{AprendizadoRL}      % ANTES do babel — fontspec precisa vir primeiro
\usepackage{babel}
\babelprovide[import, main]{brazilian}

\title[Aula 3 — Avaliação sem modelo]{Aula 3: Avaliação de Política Sem Modelo\\
  Estimando $\Vpi$ sem saber como o mundo funciona}
\subtitle[Aprendizagem por Reforço]{Monte Carlo \textbullet\ Diferença Temporal \textbullet\ Equivalência de Certeza}
\author[Prof. Marcelo K. Albertini]{Prof. Marcelo Keese Albertini \\
  \footnotesize Universidade Federal de Uberlândia \\
  \footnotesize adaptado de Emma Brunskill, CS234 (Stanford)}
\institute{Aprendizagem por Reforço}
\date{2026}

% Nota discreta: as perguntas ficam para o professor conduzir em aula
\newcommand{\paradiscussao}{\vspace{-0.2em}
  {\footnotesize\color{rlCinza}\itshape Pergunta para discussão conduzida em aula.}}

\begin{document}

\begin{frame}[plain, noframenumbering]\titlepage\end{frame}

%% ------------------------------------------------------------
\begin{frame}{Roteiro}
  \begin{itemize}
    \item \textbf{Aula passada:} processos de recompensa e de decisão de Markov;
      avaliação e controle de políticas \emph{com} o modelo verdadeiro
      (dinâmica $P$ e recompensa $\rec$ conhecidas).
    \item \textbf{Hoje:} avaliar uma política $\pol$ \alert{sem} conhecer
      $P$ nem $\rec$ — apenas interagindo com o ambiente.
    \item \textbf{Próxima aula:} \emph{controle} sem modelo — encontrar uma boa
      política a partir da experiência.
  \end{itemize}

  \begin{ideiabox}
    Até agora, o agente ``lia o manual do mundo''. A partir de hoje,
    ele precisa aprender \emph{vivendo}: executa $\pol$, observa estados,
    ações e recompensas, e estima o quanto essa política vale.
  \end{ideiabox}

  \medskip
  {\small\color{rlCinza} Leitura recomendada: Sutton \& Barto,
  seções 5.1, 5.5 e 6.1–6.3; David Silver, aula 4 (\emph{Model-Free Prediction}).}
\end{frame}

%% ------------------------------------------------------------
\begin{frame}{Avaliação por experiência direta}
  \textbf{Problema.} Estimar o retorno esperado de $\pol$ usando apenas
  dados coletados do ambiente (experiência direta).
  \begin{itemize}\setlength{\itemsep}{0.2ex}
    \item Em quase toda aplicação real (robótica, saúde, recomendação),
      \alert{não temos} o modelo verdadeiro do mundo.
    \item Hoje os dados vêm da própria $\pol$ (\emph{on-policy});
      o caso \emph{off-policy} fica para adiante.
  \end{itemize}

  \begin{quizbox}{}
    Que propriedades um bom método de avaliação de política deveria ter?
  \end{quizbox}
  \paradiscussao

  \pause
  \begin{respbox}
    Consistência, eficiência estatística, atualização barata e pouca
    memória — a régua que usaremos na aula inteira.
  \end{respbox}
\end{frame}

%% ------------------------------------------------------------
\begin{frame}{O cardápio de hoje}
  Quatro maneiras de estimar $\Vpi$ sem os modelos verdadeiros:
  \begin{enumerate}
    \item \textbf{Monte Carlo (MC):} valor $=$ média dos retornos observados.
    \item \textbf{Diferença Temporal (TD):} atualiza a cada passo,
      usando a própria estimativa como ``palpite'' do futuro.
    \item \textbf{Equivalência de certeza:} estima um modelo por máxima
      verossimilhança e planeja nele como se fosse o verdadeiro.
    \item \textbf{Avaliação em lote} (\emph{batch}): o que muda quando os
      dados são fixos e finitos — e um exemplo em que MC e TD discordam.
  \end{enumerate}

  \begin{ideiabox}
    Os três primeiros ocupam cantos diferentes de um mesmo mapa:
    quanto do futuro cada atualização enxerga (profundidade) e
    quantos futuros ela considera (largura). Voltaremos a esse mapa
    no fim da aula.
  \end{ideiabox}
\end{frame}

%% ============================================================
\section{Recordando a aula anterior}

\begin{frame}{Aquecimento: iteração de valor $\times$ iteração de política}
  \begin{quizbox}{}
    Em um MDP tabular, assintoticamente, a iteração de valor sempre produz
    uma política com o \textbf{mesmo valor} da política devolvida pela
    iteração de política. Verdadeiro ou falso?
  \end{quizbox}
  \paradiscussao

  \pause
  \begin{respbox}
    \textbf{Verdadeiro.} Ambas convergem para a função valor ótima $\Vstar$
    e devolvem uma política com valor ótimo — os \emph{caminhos} até lá é
    que diferem.
  \end{respbox}
\end{frame}

%% ------------------------------------------------------------
\begin{frame}{Aquecimento: quantas iterações a iteração de valor pode levar?}
  \begin{quizbox}{}
    A iteração de valor pode precisar de \textbf{mais} de
    $\abs{\gA}^{\abs{\gS}}$ iterações para computar a função valor ótima?
    (Suponha $\abs{\gA}$ e $\abs{\gS}$ pequenos o bastante para cada rodada
    ser exata.)
  \end{quizbox}
  \paradiscussao

  \pause
  \begin{respbox}
    \textbf{Sim.} Exemplo: MDP com um único estado e uma única ação,
    $r(s,a) = 1$, $\gamma = 0{,}9$ e $V_0(s) = 0$. Aqui
    $\abs{\gA}^{\abs{\gS}} = 1$, mas $\Vstar(s) = \tfrac{1}{1-\gamma} = 10$,
    e após a primeira iteração $V_1(s) = 1$: a convergência ao valor é
    \emph{assintótica}. (A iteração de \emph{política}, ao contrário, para
    em no máximo $\abs{\gA}^{\abs{\gS}}$ passos.)
  \end{respbox}
\end{frame}

%% ------------------------------------------------------------
\begin{frame}{Retorno e funções valor}
  \begin{defbox}{Retorno $G_t$}
    Soma descontada das recompensas do passo $t$ em diante:
    $G_t = r_t + \gamma r_{t+1} + \gamma^2 r_{t+2} + \gamma^3 r_{t+3} + \cdots$
  \end{defbox}

  \begin{defbox}{Função valor de estado $\Vpi(s)$}
    Retorno esperado ao partir de $s$ e seguir $\pol$:
    \[ \Vpi(s) = \E_{\pol}\left[\, G_t \mid s_t = s \,\right] \]
  \end{defbox}

  \begin{defbox}{Função valor de estado-ação $\Qpi(s,a)$}
    Retorno esperado ao partir de $s$, executar $a$ e depois seguir $\pol$:
    $\Qpi(s,a) = \E_{\pol}\left[\, G_t \mid s_t = s,\, a_t = a \,\right]$
  \end{defbox}
\end{frame}

%% ------------------------------------------------------------
\begin{frame}{Programação dinâmica para avaliação (com modelo)}
  Se conhecemos $P$ e $\rec$, a equação de Bellman dá um algoritmo iterativo:
  \[
    V^{\pol}_{k}(s) \;=\; \termo{\rec(s,\pol(s))}{recompensa imediata}
      \;+\; \gamma \termo{\sum_{s' \in \gS} \tran{s'}{s,\pol(s)}\, V^{\pol}_{k-1}(s')}{estimativa do futuro}
  \]

  \begin{itemize}
    \item $V^{\pol}_k(s)$ é exatamente o valor de $s$ com horizonte $k$;
      antes da convergência, é uma \emph{estimativa} de $\Vpi(s)$.
    \item O termo da direita substitui a esperança do retorno futuro
      $\E_{\pol}[r_{t+1} + \gamma r_{t+2} + \cdots]$ pela estimativa
      $V^{\pol}_{k-1}$ da iteração anterior.
  \end{itemize}

  \begin{defbox}{\emph{Bootstrapping}}
    Atualizar uma estimativa usando \emph{outra estimativa} no lugar da
    quantidade verdadeira. A PD faz bootstrapping; guarde o termo — ele
    é a alma do TD.
  \end{defbox}
\end{frame}

%% ------------------------------------------------------------
\begin{frame}{O que a PD usa — e o que não teremos hoje}
  \centering
  \backupdiagrama

  \medskip
  \begin{itemize}
    \item A PD calcula a média \textbf{sobre todos} os próximos estados $s'$,
      ponderada por $\tran{s'}{s,\pol(s)}$ — para isso, precisa do modelo.
    \item Hoje só temos \alert{amostras}: um $s'$ por vez, sorteado pelo
      próprio ambiente.
  \end{itemize}

  \begin{alertabox}
    Sem o modelo, não dá para somar sobre $s'$ — não sabemos as
    probabilidades. Toda a aula é sobre como contornar isso.
  \end{alertabox}
\end{frame}

%% ============================================================
\section{Avaliação de política por Monte Carlo}

\begin{frame}{A ideia mais simples possível}
  \[
    \Vpi(s) = \E_{\tau \sim \pol}\left[\, G_t \mid s_t = s \,\right]
    \qquad \text{(esperança sobre trajetórias $\tau$ geradas por $\pol$)}
  \]

  \begin{ideiabox}
    Valor é uma esperança — e esperanças se estimam com médias.
    Rode a política várias vezes, anote o retorno obtido a partir de cada
    estado e tire a média: \alert{valor $=$ retorno médio}.
  \end{ideiabox}

  \begin{itemize}
    \item Não requer modelo de dinâmica nem de recompensa.
    \item Nem sequer assume que o estado é Markoviano!
    \item Requer episódios \textbf{finitos}: só podemos calcular $G_t$
      quando o episódio termina.
    \item Episódios podem ter comprimentos diferentes (estados terminais).
  \end{itemize}
\end{frame}

%% ------------------------------------------------------------
\begin{frame}{MC de primeira visita}
  \begin{algbox}{MC de primeira visita (\emph{first-visit})}
    \begin{itemize}\setlength{\itemsep}{0.2ex}
      \item Inicialize $N(s) = 0$, $G(s) = 0$ para todo $s \in \gS$
      \item Repita:
      \begin{itemize}\setlength{\itemsep}{0.15ex}
        \item Amostre um episódio $i$: $s_{i,1}, a_{i,1}, r_{i,1}, s_{i,2}, \ldots, s_{i,T_i}, a_{i,T_i}, r_{i,T_i}$
        \item Calcule $G_{i,t} = r_{i,t} + \gamma r_{i,t+1} + \gamma^2 r_{i,t+2} + \cdots + \gamma^{T_i - t} r_{i,T_i}$
        \item Para cada passo $t$ do episódio:
          se $t$ é a \alert{primeira} visita ao estado $s$ neste episódio:
          \begin{itemize}
            \item $N(s) \leftarrow N(s) + 1$; \quad $G(s) \leftarrow G(s) + G_{i,t}$
            \item $\Vpi(s) \leftarrow G(s)/N(s)$
          \end{itemize}
      \end{itemize}
    \end{itemize}
  \end{algbox}

  Cada estado contribui \textbf{no máximo uma amostra por episódio} —
  as amostras usadas na média são independentes entre episódios.
\end{frame}

%% ------------------------------------------------------------
\begin{frame}{MC de toda visita}
  \begin{algbox}{MC de toda visita (\emph{every-visit})}
    \begin{itemize}\setlength{\itemsep}{0.2ex}
      \item Idêntico ao anterior, exceto pela condição destacada:
      \item Para cada passo $t$ do episódio, com $s = s_{i,t}$
        (\alert{toda} visita conta, não só a primeira):
      \begin{itemize}
        \item $N(s) \leftarrow N(s) + 1$; \quad $G(s) \leftarrow G(s) + G_{i,t}$
        \item $\Vpi(s) \leftarrow G(s)/N(s)$
      \end{itemize}
    \end{itemize}
  \end{algbox}

  \begin{itemize}
    \item Se a política revisita $s$ dentro do mesmo episódio, cada visita
      gera uma amostra de retorno.
    \item Essas amostras são \emph{correlacionadas} (compartilham o fim do
      episódio) — isso terá consequências estatísticas, como veremos.
  \end{itemize}
\end{frame}

%% ------------------------------------------------------------
\begin{frame}{Exemplo resolvido: \emph{Mars Rover}}
  \centering
  \scalebox{0.9}{\marsrover{1,7}}

  \medskip
  \begin{exbox}{— MC no Mars Rover}
    Recompensas $\rec = [+1,\ 0,\ 0,\ 0,\ 0,\ 0,\ +10]$ para qualquer ação;
    $s_1$ e $s_7$ terminam o episódio. Trajetória observada:
    \[
      (s_3, a_1, 0,\ s_2, a_1, 0,\ s_2, a_1, 0,\ s_1, a_1, +1,\ \text{terminal})
    \]
    Com $\gamma < 1$, calcule as estimativas MC de \textbf{primeira visita}
    e de \textbf{toda visita} para $V(s_2)$.
  \end{exbox}
\end{frame}

%% ------------------------------------------------------------
\begin{frame}{Exemplo resolvido: solução}
  Retornos a partir de cada passo da trajetória
  $(s_3, 0)\!\to\!(s_2, 0)\!\to\!(s_2, 0)\!\to\!(s_1, +1)\!\to\!\text{fim}$:
  \begin{align*}
    G_{1} &= 0 + \gamma\cdot 0 + \gamma^2\cdot 0 + \gamma^3\cdot 1 = \gamma^3
      && \text{(a partir de $s_3$)}\\
    G_{2} &= 0 + \gamma\cdot 0 + \gamma^2\cdot 1 = \gamma^2
      && \text{(1ª visita a $s_2$)}\\
    G_{3} &= 0 + \gamma\cdot 1 = \gamma
      && \text{(2ª visita a $s_2$)}\\
    G_{4} &= 1 && \text{(a partir de $s_1$)}
  \end{align*}

  \begin{itemize}
    \item \textbf{Primeira visita:} só $G_2$ conta para $s_2$:\quad
      $\mdestaque{rlLaranja}{V(s_2) = \gamma^2}$
    \item \textbf{Toda visita:} média de $G_2$ e $G_3$:\quad
      $\mdestaque{rlLaranja}{V(s_2) = \tfrac{\gamma^2 + \gamma}{2}}$
  \end{itemize}

  \medskip
  Além disso: $V(s_3) = \gamma^3$, $V(s_1) = 1$, e os estados nunca
  visitados ($s_4$ a $s_7$) permanecem na inicialização.
\end{frame}

%% ------------------------------------------------------------
\begin{frame}{MC incremental: a mesma média, sem guardar somas}
  A média pode ser atualizada amostra a amostra. Ao visitar $s$ e observar
  o retorno $G_{i,t}$, com $N(s) \leftarrow N(s)+1$:
  \[
    \Vpi(s) \;=\; \Vpi(s)\,\frac{N(s)-1}{N(s)} + \frac{G_{i,t}}{N(s)}
    \;=\; \Vpi(s) + \frac{1}{N(s)}\bigl(G_{i,t} - \Vpi(s)\bigr)
  \]

  Generalizando $\tfrac{1}{N(s)}$ para uma taxa de aprendizado $\alpha$:
  \[
    \mdestaque{rlLaranja}{\Vpi(s_{i,t}) \;=\; \Vpi(s_{i,t})
      + \alpha\,\bigl(G_{i,t} - \Vpi(s_{i,t})\bigr)}
  \]

  \begin{ideiabox}
    Este é o molde de quase tudo que vem pela frente no curso:
    $\;\text{estimativa} \leftarrow \text{estimativa}
      + \alpha\,(\text{alvo} - \text{estimativa})$,
    com uma \emph{taxa de aprendizado} $\alpha$ e um \emph{alvo}.
    Muda-se o alvo, nasce um algoritmo novo.
  \end{ideiabox}
\end{frame}

%% ------------------------------------------------------------
\begin{frame}{Interlúdio estatístico: como julgar um estimador?}
  Modelo estatístico com parâmetro $\theta$ gera dados $x \sim P(x \mid \theta)$;
  um estimador $\hat{\theta} = \hat{\theta}(x)$ tenta recuperar $\theta$.

  \begin{defbox}{Viés, variância e EQM}
    \vspace{-1.2em}
    \begin{align*}
      \text{Vi\'es}(\hat{\theta})
        &= \E_{x\mid\theta}[\hat{\theta}\,] - \theta\\
      \operatorname{Var}(\hat{\theta})
        &= \E_{x\mid\theta}\bigl[(\hat{\theta} - \E[\hat{\theta}\,])^2\bigr]\\
      \operatorname{EQM}(\hat{\theta})
        &= \operatorname{Var}(\hat{\theta})
          + \text{Vi\'es}(\hat{\theta})^2
    \end{align*}
  \end{defbox}

  \begin{defbox}{Estimador consistente}
    Com $n$ amostras e estimativa $\hat{\theta}_n$: para todo
    $\epsilon > 0$,\;
    $\lim_{n \to \infty} \Prob\bigl(\abs{\hat{\theta}_n - \theta} > \epsilon\bigr) = 0$.
  \end{defbox}
\end{frame}

%% ------------------------------------------------------------
\begin{frame}{Viés zero implica consistência?}
  \begin{quizbox}{}
    Se um estimador é não-viesado (viés $=0$), ele é necessariamente
    consistente?
  \end{quizbox}
  \paradiscussao

  \pause
  \begin{respbox}
    \textbf{Não.} Exemplo: estimar a média de uma Gaussiana usando
    \emph{apenas a primeira} observação, $\hat{\theta}_n = x_1$.
    É não-viesado ($\E[x_1] = \theta$), mas ignora os demais dados —
    a distribuição do erro não encolhe quando $n \to \infty$.
    Consistência exige que o estimador \emph{aproveite} os dados novos.
  \end{respbox}

  \begin{ideiabox}
    Viés e variância medem coisas diferentes: viés é ``mirar no lugar
    errado''; variância é ``tremer a mão''. O EQM cobra os dois.
  \end{ideiabox}
\end{frame}

%% ------------------------------------------------------------
\begin{frame}{Propriedades dos estimadores MC}
  \begin{itemize}
    \item \textbf{Primeira visita:} estimador \alert{não-viesado} de
      $\E_{\pol}[G_t \mid s_t = s]$. Pela lei dos grandes números,
      $\Vpi(s) \to \E_{\pol}[G_t \mid s_t = s]$ quando $N(s) \to \infty$.
    \item \textbf{Toda visita:} \alert{viesado} (amostras correlacionadas
      dentro do episódio), porém \textbf{consistente} — e com frequência
      tem EQM \emph{menor} na prática.
    \item \textbf{Incremental:} as propriedades dependem da taxa $\alpha$.
  \end{itemize}

  \begin{teobox}{Condições de convergência (Robbins–Monro)}
    Permitindo $\alpha$ variar por atualização
    ($\alpha_n(s_j)$ na $n$-ésima atualização do estado $s_j$),
    o MC incremental converge para $\Vpi(s_j)$ se
    \[
      \sum_{n=1}^{\infty} \alpha_n(s_j) = \infty
      \qquad \text{e} \qquad
      \sum_{n=1}^{\infty} \alpha_n^2(s_j) < \infty.
    \]
  \end{teobox}
\end{frame}

%% ------------------------------------------------------------
\begin{frame}{Limitações do Monte Carlo}
  \begin{alertabox}
    O retorno $G_t$ acumula a aleatoriedade de \emph{todas} as ações,
    transições e recompensas até o fim do episódio: o estimador MC tem,
    em geral, \textbf{variância alta}.
  \end{alertabox}

  \begin{itemize}
    \item Reduzir essa variância pode exigir \emph{muitos} dados —
      impraticável quando coletar dados é caro, lento ou arriscado
      (medicina, robôs físicos, educação).
    \item Exige cenário \textbf{episódico}: o episódio precisa
      \emph{terminar} antes de qualquer atualização.
    \item Nada de tarefas contínuas / horizonte infinito.
  \end{itemize}
\end{frame}

%% ------------------------------------------------------------
\begin{frame}{Resumo: Monte Carlo}
  \begin{itemize}
    \item Estima a esperança $\Vpi(s) = \E_{\pol}[G_t \mid s_t = s]$
      por média empírica de retornos, com episódios amostrados de $\pol$.
    \item Simples; não requer modelo; \textbf{não assume Markov}.
    \item Converge para o valor verdadeiro sob hipóteses brandas.
    \item Às vezes é preferível à programação dinâmica \emph{mesmo com o
      modelo em mãos} — por exemplo, quando $\gS$ é gigantesco e só nos
      interessam alguns estados, ou quando o estado não é Markoviano.
  \end{itemize}
\end{frame}

%% ============================================================
\section{Aprendizagem por Diferença Temporal}

\begin{frame}{A ideia central do RL}
  Sutton e Barto apontam a aprendizagem por diferença temporal (TD) como
  \emph{a} ideia mais central e original de toda a aprendizagem por reforço.

  \medskip
  \begin{itemize}
    \item Combinação de Monte Carlo (\textbf{amostra} do ambiente) com
      programação dinâmica (\textbf{bootstrapping}).
    \item Livre de modelo.
    \item Funciona em cenários episódicos \emph{e} de horizonte infinito.
    \item Atualiza $V$ \alert{imediatamente} após cada tupla $(s, a, r, s')$
      — não espera o episódio acabar.
  \end{itemize}

  \begin{ideiabox}
    O MC espera o fim do episódio para saber ``quanto valeu''.
    O TD não espera: usa a recompensa de um passo e o \emph{próprio
    palpite} sobre o estado seguinte como substituto do resto do futuro.
  \end{ideiabox}
\end{frame}

%% ------------------------------------------------------------
\begin{frame}{Do MC ao TD em uma linha}
  Atualização do MC incremental (alvo $=$ retorno completo $G_{i,t}$):
  \[
    \Vpi(s) = \Vpi(s) + \alpha\,\bigl(\textcolor{rlCinza}{G_{i,t}} - \Vpi(s)\bigr)
  \]

  Mas $G_t = r_t + \gamma G_{t+1}$, e $\E[G_{t+1}] $ é o valor do próximo
  estado... que já estamos estimando! Troque o futuro observado pela
  estimativa atual:
  \[
    \Vpi(s_t) = \Vpi(s_t)
      + \alpha\,\bigl(\mdestaque{rlLaranja}{r_t + \gamma \Vpi(s_{t+1})} - \Vpi(s_t)\bigr)
  \]

  \begin{defbox}{Alvo TD e erro TD}
    O \textbf{alvo TD} é $r_t + \gamma \Vpi(s_{t+1})$; o \textbf{erro TD} é
    \[
      \delta_t = r_t + \gamma \Vpi(s_{t+1}) - \Vpi(s_t).
    \]
    A atualização é simplesmente $\Vpi(s_t) \leftarrow \Vpi(s_t) + \alpha\,\delta_t$.
  \end{defbox}
\end{frame}

%% ------------------------------------------------------------
\begin{frame}{Algoritmo TD(0)}
  \begin{algbox}{TD(0), ou TD de 1 passo}
    \begin{itemize}\setlength{\itemsep}{0.2ex}
      \item Entrada: taxa de aprendizado $\alpha$
      \item Inicialize $\Vpi(s) = 0$ para todo $s \in \gS$
      \item Repita:
      \begin{itemize}
        \item Amostre uma tupla $(s_t, a_t, r_t, s_{t+1})$ executando $\pol$
        \item $\Vpi(s_t) \leftarrow \Vpi(s_t)
          + \alpha\bigl(\,\termo{r_t + \gamma \Vpi(s_{t+1})}{alvo TD} - \Vpi(s_t)\bigr)$
      \end{itemize}
    \end{itemize}
  \end{algbox}

  \begin{itemize}
    \item Cada tupla é usada uma vez, na hora — não é preciso armazenar
      episódios nem esperar que terminem.
    \item Por isso serve para tarefas contínuas, sem estados terminais.
    \item O ``(0)'' indica olhar 1 passo à frente; há uma família inteira,
      TD($\lambda$), interpolando entre TD(0) e Monte Carlo.
  \end{itemize}
\end{frame}

%% ------------------------------------------------------------
\begin{frame}{Exemplo resolvido: TD no Mars Rover}
  Mesmo cenário: $\rec = [+1, 0, 0, 0, 0, 0, +10]$,
  $\pol(s) = a_1$ para todo $s$, e $s_1, s_7$ terminam o episódio.
  Mesma trajetória, processada agora tupla a tupla, com $\alpha = 1$
  e $V$ inicializado em zero:

  \medskip
  \centering
  \begin{tikzpicture}[node distance=9mm]
    \node[estado] (a) {$s_3$};
    \node[estado, right=of a] (b) {$s_2$};
    \node[estado, right=of b] (c) {$s_2$};
    \node[estado destacado, right=of c] (d) {$s_1$};
    \node[estado terminal, right=of d] (e) {$\top$};
    \draw[trans] (a) -- node[probl, above=1pt] {$r=0$} (b);
    \draw[trans] (b) -- node[probl, above=1pt] {$r=0$} (c);
    \draw[trans] (c) -- node[probl, above=1pt] {$r=0$} (d);
    \draw[trans destac] (d) -- node[recomp, above=2pt] {$r=+1$} (e);
  \end{tikzpicture}

  \medskip
  \raggedright
  \begin{enumerate}\setlength{\itemsep}{0.2ex}
    \item $(s_3, 0, s_2)$:\quad $V(s_3) \leftarrow 0 + \gamma\cdot 0 = 0$
    \item $(s_2, 0, s_2)$:\quad $V(s_2) \leftarrow 0 + \gamma\cdot 0 = 0$
    \item $(s_2, 0, s_1)$:\quad $V(s_2) \leftarrow 0 + \gamma\cdot 0 = 0$
    \item $(s_1, +1, \top)$:\quad $V(s_1) \leftarrow 1 + \gamma\cdot 0 = 1$
  \end{enumerate}

  Resultado: $V = [1,\ 0,\ 0,\ 0,\ 0,\ 0,\ 0]$.
\end{frame}

%% ------------------------------------------------------------
\begin{frame}{MC e TD divergem — no mesmo episódio!}
  Para a mesma trajetória (tomando $\gamma = 1$ para simplificar):
  \begin{center}
  \renewcommand{\arraystretch}{1.3}
  \begin{tabular}{lccc}
    & $V(s_1)$ & $V(s_2)$ & $V(s_3)$ \\ \hline
    MC (primeira visita) & $1$ & $1$ & $1$ \\
    TD(0), $\alpha = 1$  & $1$ & $0$ & $0$ \\
  \end{tabular}
  \end{center}

  \begin{itemize}
    \item O MC propaga o desfecho $+1$ \textbf{para trás}, a todos os
      estados visitados no episódio.
    \item O TD(0) usa cada tupla \textbf{uma única vez}: quando $s_2$ foi
      atualizado, $V(s_1)$ ainda era $0$ — a notícia não teve tempo de
      se propagar.
    \item Com mais episódios (ou repassando os dados), o TD também propaga
      a informação; a diferença é a \emph{velocidade com que cada dado
      influencia} a estimativa.
  \end{itemize}
\end{frame}

%% ------------------------------------------------------------
\begin{frame}{Teste sua compreensão: o papel de $\alpha$}
  \begin{quizbox}{}
    Quais afirmações sobre o TD(0) são verdadeiras?
    \begin{enumerate}\setlength{\itemsep}{0ex}
      \item Com $\alpha = 0$, o alvo TD pesa mais que a estimativa antiga.
      \item Com $\alpha = 1$, a estimativa é substituída pelo alvo TD.
      \item Com $\alpha = 1$, em estados com vários sucessores possíveis,
        $V$ pode oscilar para sempre.
      \item Há MDPs determinísticos em que o TD com $\alpha = 1$ converge.
    \end{enumerate}
  \end{quizbox}
  \paradiscussao

  \pause
  \begin{respbox}
    \textbf{2, 3 e 4.} Com $\alpha = 0$ nada é atualizado (1 é falsa);
    com $\alpha = 1$ a estimativa vira o alvo (2); alvo estocástico faz
    $V$ quicar (3); alvo determinístico é fixo e converge (4).
  \end{respbox}
\end{frame}

%% ------------------------------------------------------------
\begin{frame}{Resumo: Diferença Temporal}
  \begin{itemize}
    \item Amostra do ambiente (como MC) \textbf{e} faz bootstrapping (como PD).
    \item Livre de modelo; serve para cenários episódicos e contínuos.
    \item Atualiza a cada tupla $(s, a, r, s')$ — aprendizado \emph{online}.
    \item Estimador \alert{viesado}: no começo, o alvo depende da
      inicialização de $V$.
    \item Em geral, variância \textbf{bem menor} que a do MC: o alvo TD
      contém uma única ação, recompensa e transição aleatórias, e não a
      sequência inteira.
    \item Consistente se $\alpha$ satisfaz as condições de Robbins–Monro
      (as mesmas do MC incremental).
  \end{itemize}
\end{frame}

%% ============================================================
\section{Equivalência de certeza}

\begin{frame}{E se estimássemos o modelo?}
  Terceira via, \emph{baseada em modelo}: aprender $\hat{P}$ e $\hat{r}$
  dos dados e \textbf{planejar} como se fossem verdadeiros.

  \begin{defbox}{Modelo de máxima verossimilhança (contagens)}
    Para cada par $(s,a)$ com $N(s,a)$ visitas até agora:
    \[
      \hat{P}(s' \mid s, a) = \frac{\#\,\{(s,a) \to s'\}}{N(s,a)},
      \qquad
      \hat{r}(s, a) = \frac{\text{soma das recompensas de } (s,a)}{N(s,a)}.
    \]
  \end{defbox}

  \begin{algbox}{Avaliação por equivalência de certeza}
    \begin{itemize}\setlength{\itemsep}{0.2ex}
      \item A cada tupla $(s, a, r, s')$: atualize contagens e modelo
        $(\hat{P}, \hat{r})$
      \item Calcule $\Vpi$ no MDP estimado com qualquer método de PD da
        aula 2 \emph{(inicializar todos os pares $(s,a)$)}
    \end{itemize}
  \end{algbox}
\end{frame}

%% ------------------------------------------------------------
\begin{frame}{Equivalência de certeza: custo e propriedades}
  \begin{itemize}
    \item \textbf{Custo alto:} replanejar a cada atualização custa
      $O(\abs{\gS}^3)$ pela solução matricial analítica, ou
      $O(\abs{\gS}^2 \abs{\gA})$ por iteração dos métodos iterativos.
    \item \textbf{Eficiência estatística excelente:} cada tupla melhora o
      modelo inteiro, e o planejamento espreme toda a informação disponível.
    \item \textbf{Consistente} quando o processo é de fato Markoviano.
    \item Estende-se com naturalidade à avaliação \emph{off-policy}:
      o modelo estimado não depende de qual política gerou os dados.
  \end{itemize}

  \begin{ideiabox}
    O nome diz tudo: tratamos a estimativa \emph{como se} tivéssemos
    certeza dela. Muito eficiente com poucos dados; muito caro em
    computação — o oposto exato do TD.
  \end{ideiabox}
\end{frame}

%% ------------------------------------------------------------
\begin{frame}{Exemplo resolvido: equivalência de certeza no Mars Rover}
  \framesubtitle{Mesma trajetória:
    $(s_3, a_1, 0,\ s_2, a_1, 0,\ s_2, a_1, 0,\ s_1, a_1, +1,\ \text{terminal})$}
  \textbf{Modelo estimado por contagens:}
  \begin{itemize}\setlength{\itemsep}{0.15ex}
    \item $\hat{r} = [+1,\ 0,\ 0,\ 0,\ 0,\ 0,\ 0]$
      (só vimos recompensas em $s_1$, $s_2$, $s_3$)
    \item $\hat{p}(s_2 \mid s_3, a_1) = 1$;\quad
      $\hat{p}(\text{terminal} \mid s_1, a_1) = 1$
    \item $s_2$: duas visitas, uma ficou e outra desceu $\Rightarrow$
      $\hat{p}(s_2 \mid s_2, a_1) = \hat{p}(s_1 \mid s_2, a_1) = \tfrac{1}{2}$
  \end{itemize}

  \textbf{Planejando no modelo} (resolvendo $V = \hat{r} + \gamma \hat{P} V$):
  \vspace{-0.4em}
  \begin{align*}
    V(s_1) &= 1\\
    V(s_2) &= \gamma\bigl(\tfrac12 V(s_2) + \tfrac12 V(s_1)\bigr)
      \;\Rightarrow\; V(s_2) = \tfrac{\gamma}{2 - \gamma}
    \qquad
    V(s_3) = \gamma\, V(s_2) = \tfrac{\gamma^2}{2-\gamma}
  \end{align*}
  \vspace{-1.2em}

  Note: \textbf{um único episódio} já deu estimativas não-nulas aos três
  estados — nenhum dos métodos anteriores conseguiu isso.
\end{frame}

%% ============================================================
\section{Avaliação em lote}

\begin{frame}{O cenário em lote (\emph{batch})}
  Até aqui, os dados chegavam em fluxo. Agora suponha um conjunto
  \textbf{fixo e finito} de $K$ episódios coletados sob $\pol$:
  \vspace{-0.5em}
  \[
    \{\, s_{k,1}, a_{k,1}, r_{k,1}, \ldots, s_{k,T_k} \,\}_{k=1}^{K}
  \]
  \vspace{-1.4em}
  \begin{itemize}\setlength{\itemsep}{0.2ex}
    \item Regime típico quando coletar dados novos é caro ou arriscado
      (prontuários médicos, logs de um sistema em produção).
    \item Ideia: \textbf{repassar} os mesmos episódios várias vezes,
      aplicando MC ou TD(0) até as estimativas estabilizarem.
  \end{itemize}
  \vspace{-0.4em}

  \begin{quizbox}{}
    Repassando os mesmos dados infinitas vezes, MC e TD(0) convergem
    para a \emph{mesma} resposta?
  \end{quizbox}
  \paradiscussao
\end{frame}

%% ------------------------------------------------------------
\begin{frame}{Um exemplo minúsculo que separa MC de TD}
  Dois estados, $A$ e $B$; sem desconto ($\gamma = 1$); \textbf{8 episódios}
  observados:

  \begin{center}
  \renewcommand{\arraystretch}{1.2}
  \begin{tabular}{cl}
    1 episódio & $A,\ 0,\ B,\ 0,\ \text{fim}$\\
    6 episódios & $B,\ 1,\ \text{fim}$\\
    1 episódio & $B,\ 0,\ \text{fim}$\\
  \end{tabular}
  \end{center}

  \begin{exbox}{— o que os dados dizem?}
    No regime em lote, quais os valores de $V(B)$ e de $V(A)$?
  \end{exbox}

  \pause
  \begin{respbox}
    $V(B) = \tfrac{6}{8} = 0{,}75$ — MC e TD concordam.
    Para $V(A)$, eles \textbf{discordam}: MC responde $0$;
    TD(0) responde $0{,}75$.
  \end{respbox}
\end{frame}

%% ------------------------------------------------------------
\begin{frame}{Por que discordam — e o que cada um ``acredita''}
  \begin{columns}[T]
    \begin{column}{0.46\textwidth}
      \small\textbf{Monte Carlo.} Só há \emph{um} retorno observado a
      partir de $A$, e ele valeu $0$. Logo $V(A) = 0$: o MC ajusta-se
      aos \alert{retornos que aconteceram}.
    \end{column}
    \begin{column}{0.50\textwidth}
      \small\textbf{TD(0).} Nos dados, $A$ sempre levou a $B$ com $r=0$;
      então $V(A) = 0 + \gamma V(B) = 0{,}75$: o TD \alert{acredita na
      estrutura de Markov} e generaliza via $B$.
    \end{column}
  \end{columns}

  \smallskip
  \begin{teobox}[fontupper=\small]{Para onde cada método converge no lote}
    \vspace{-0.2em}
    \begin{itemize}\setlength{\itemsep}{0.1ex}
      \item \textbf{MC:} valores que minimizam o erro quadrático médio
        sobre os retornos observados.
      \item \textbf{TD(0):} $\Vpi$ do MDP de máxima verossimilhança dos
        dados — a solução da \alert{equivalência de certeza}.
    \end{itemize}
  \end{teobox}

  \begin{ideiabox}
    Se o mundo é mesmo Markoviano, a resposta do TD aproveita melhor os
    dados; se não é, a do MC é a mais honesta.
  \end{ideiabox}
\end{frame}

%% ------------------------------------------------------------
\begin{frame}{O modelo que o TD constrói sem querer}
  \framesubtitle{MDP de máxima verossimilhança estimado a partir dos 8 episódios}
  \centering
  \begin{tikzpicture}[node distance=22mm]
    \node[estado] (A) {$A$};
    \node[estado, right=of A] (B) {$B$};
    \node[estado terminal, right=of B] (T) {$\top$};
    \draw[trans destac] (A) -- node[probl, above=1pt] {$\hat p = 1$}
                             node[recomp, below=3pt] {$\hat r = 0$} (B);
    \draw[trans] (B) -- node[probl, above=1pt] {$\hat p = 1$}
                       node[recomp, below=3pt] {$\hat r = 0{,}75$} (T);
  \end{tikzpicture}

  \smallskip
  \raggedright
  \begin{itemize}\setlength{\itemsep}{0.2ex}
    \item Das 8 trajetórias, a única que passou por $A$ foi para $B$
      com recompensa $0$ $\Rightarrow$ $\hat p(B \mid A) = 1$, $\hat r(A) = 0$.
    \item De $B$, seis vezes a recompensa foi $1$ e duas vezes foi $0$
      $\Rightarrow$ $\hat r(B) = 6/8 = 0{,}75$.
    \item Planejar nesse modelo dá $V(A) = 0 + \gamma\,V(B) = 0{,}75$ —
      \alert{exatamente a resposta do TD(0) em lote}.
  \end{itemize}

  \begin{ideiabox}
    O TD nunca escreve $\hat P$ nem $\hat r$ na memória, mas converge para
    a mesma resposta que a equivalência de certeza daria. É o preço e o
    prêmio de assumir Markov.
  \end{ideiabox}
\end{frame}

%% ============================================================
\section{Comparando os métodos}

\begin{frame}{Viés e variância dos alvos}
  \begin{center}
  \renewcommand{\arraystretch}{1.35}
  \begin{tabular}{lcc}
    \textbf{Alvo} & \textbf{Viés} & \textbf{Variância}\\ \hline
    Retorno $G_t$ (MC)
      & nenhum (1ª visita) & alta: episódio inteiro é aleatório\\
    $r_t + \gamma \Vpi(s_{t+1})$ (TD)
      & sim (depende de $\Vpi$ atual) & baixa: 1 passo aleatório\\
  \end{tabular}
  \end{center}

  \begin{itemize}
    \item \textbf{MC:} não-viesado (primeira visita), variância alta,
      consistente — \emph{inclusive} com aproximação de função.
    \item \textbf{TD(0):} viesado, variância menor; converge para $\Vpi$
      na representação tabular, mas \alert{pode não convergir com
      aproximação de função} (tema de aulas futuras).
  \end{itemize}

  \begin{ideiabox}
    É o dilema viés–variância clássico, vestido de RL: o TD aceita um
    pouco de viés em troca de aprender mais rápido e com menos ruído.
  \end{ideiabox}
\end{frame}

%% ------------------------------------------------------------
\begin{frame}{Um mapa para os três métodos}
  \framesubtitle{Quantos futuros cada atualização considera (largura)
    $\times$ até onde ela olha antes de usar uma estimativa (profundidade)}
  \centering
  \scalebox{0.82}{%
  \begin{tikzpicture}[x=1cm, y=1cm]
    % eixos
    \draw[trans, rlAzul] (0,0) -- (8.6,0);
    \draw[trans, rlAzul] (0,0) -- (0,4.4);
    \node[font=\scriptsize\itshape, text=rlAzul, below] at (4.25,-0.5)
      {largura do backup};
    \node[font=\scriptsize\itshape, text=rlAzul, rotate=90, above] at (-0.55,2.2)
      {profundidade do backup};
    \node[font=\scriptsize, text=rlCinza, below] at (1.9,-0.03) {uma amostra};
    \node[font=\scriptsize, text=rlCinza, below] at (6.6,-0.03) {esperança completa};
    \node[font=\scriptsize, text=rlCinza, rotate=90, above] at (-0.03,1.1) {1 passo};
    \node[font=\scriptsize, text=rlCinza, rotate=90, above] at (-0.03,3.6) {episódio};
    % quadrantes
    \node[draw=rlLaranja, very thick, fill=rlLaranja!12, rounded corners=1.5mm,
          minimum width=2.5cm, minimum height=9mm, font=\small\bfseries]
          at (1.9,1.1) {TD(0)};
    \node[draw=rlAzulClaro, very thick, fill=rlAzulClaro!10, rounded corners=1.5mm,
          minimum width=3.2cm, minimum height=9mm, font=\small\bfseries]
          at (6.6,1.1) {PD (com modelo)};
    \node[draw=rlTeal, very thick, fill=rlTeal!10, rounded corners=1.5mm,
          minimum width=2.5cm, minimum height=9mm, font=\small\bfseries]
          at (1.9,3.6) {Monte Carlo};
    \node[draw=rlCinza, thick, dashed, fill=rlFundo, rounded corners=1.5mm,
          minimum width=3.2cm, minimum height=9mm, font=\small]
          at (6.6,3.6) {busca exaustiva};
  \end{tikzpicture}}

  \smallskip
  \raggedright
  A equivalência de certeza vive perto da PD: estima a esperança completa
  — mas num modelo aprendido dos dados.
\end{frame}

%% ------------------------------------------------------------
\begin{frame}{Tabela comparativa}
  \centering
  \renewcommand{\arraystretch}{1.15}
  {\footnotesize
  \begin{tabular}{lccc}
    & \textbf{MC} & \textbf{TD(0)} & \textbf{Equiv. certeza}\\ \hline
    Precisa de modelo?            & não & não & não (estima um)\\
    Faz bootstrapping?            & não & sim & sim (via PD)\\
    Exige episódios finitos?      & sim & não & não\\
    Assume Markov?                & não & sim & sim\\
    Viés                          & zero (1ª visita) & sim & sim (com poucos dados)\\
    Variância                     & alta & baixa & baixa\\
    Custo por atualização         & baixo & baixíssimo & alto\\
    Eficiência estatística        & baixa & média & alta\\
    Avaliação off-policy direta?  & não & não & sim\\
  \end{tabular}}

  \smallskip
  \begin{alertabox}
    Não existe vencedor absoluto: a escolha depende de qual recurso é
    escasso — dados, computação ou memória — e de o estado ser Markoviano.
  \end{alertabox}
\end{frame}

%% ------------------------------------------------------------
\begin{frame}{O que levamos desta aula}
  \begin{itemize}
    \item Avaliação de política \textbf{sem modelo}: substituímos esperanças
      por \emph{amostras} (MC), por \emph{amostra + bootstrapping} (TD),
      ou por um \emph{modelo estimado} (equivalência de certeza).
    \item O molde universal
      $\;\text{estimativa} \leftarrow \text{estimativa}
        + \alpha(\text{alvo} - \text{estimativa})\;$
      reaparecerá na aula que vem — trocando $V$ por $Q$ e avaliação por
      \textbf{controle}.
    \item Viés, variância, consistência e custo computacional são a régua
      para comparar métodos — memorize a tabela.
  \end{itemize}

  \medskip
  \textbf{Próxima aula:} controle sem modelo — SARSA e Q-learning.

  \medskip
  {\small\color{rlCinza} Material adaptado de Emma Brunskill,
  \emph{CS234: Reinforcement Learning} (Stanford). Estrutura relacionada à
  aula 4 de David Silver; leituras em Sutton \& Barto,
  seções 5.1, 5.5 e 6.1–6.3.}
\end{frame}

%% ============================================================
\section*{Exercícios opcionais}
\begin{frame}[plain, noframenumbering]
  \begin{tikzpicture}[remember picture, overlay]
    \fill[rlAzul] (current page.south west) rectangle (current page.north east);
    \node[anchor=west, text width=0.72\paperwidth]
      at ([xshift=1.1cm]current page.west) {%
        {\color{rlAmbar}\small\bfseries MATERIAL EXTRA}\\[0.55em]
        {\color{white}\LARGE\bfseries Exercícios opcionais}\\[0.55em]
        {\color{rlLaranja}\rule{3em}{2.4pt}}};
  \end{tikzpicture}
\end{frame}

\begin{frame}{Exercício: MC incremental e suas variantes}
  \begin{exbox}{— verdadeiro ou falso}
    Considere o MC incremental
    $\Vpi(s_{i,t}) \leftarrow \Vpi(s_{i,t}) + \alpha\,(G_{i,t} - \Vpi(s_{i,t}))$.
    \begin{enumerate}\setlength{\itemsep}{0.15ex}
      \item Com $\alpha = 1$, o MC incremental equivale ao MC de primeira visita.
      \item Com $\alpha = \tfrac{1}{N(s_{i,t})}$, equivale ao MC de toda visita.
      \item Usar $\alpha > \tfrac{1}{N(s_{i,t})}$ pode ajudar em domínios
        não-estacionários.
    \end{enumerate}
  \end{exbox}

  \pause
  \begin{respbox}
    \textbf{(1) Falso.} Com $\alpha = 1$ a estimativa vira o \emph{último}
    retorno observado — não a média, e todas as visitas atualizam.
    \textbf{(2) Verdadeiro.} É exatamente a média incremental sobre todas
    as visitas.
    \textbf{(3) Verdadeiro.} Pesos maiores para dados recentes permitem
    ``esquecer'' um passado que já não descreve o ambiente.
  \end{respbox}
\end{frame}

\begin{frame}{Exercício: complete o quadro do Mars Rover}
  \begin{exbox}{— os três métodos lado a lado}
    Para a trajetória
    $(s_3, a_1, 0,\ s_2, a_1, 0,\ s_2, a_1, 0,\ s_1, a_1, +1,\ \text{terminal})$,
    com $V$ inicializado em zero:
    \begin{enumerate}\setlength{\itemsep}{0.15ex}
      \item Verifique as estimativas de primeira visita com $\gamma = 1$:
        $V = [1, 1, 1, 0, 0, 0, 0]$.
      \item Verifique o TD(0) com $\alpha = 1$:
        $V = [1, 0, 0, 0, 0, 0, 0]$.
      \item Refaça a conta da equivalência de certeza resolvendo
        $V = (I - \gamma \hat{P})^{-1} \hat{r}$ restrita aos estados
        $\{s_1, s_2, s_3, \text{terminal}\}$ e confirme
        $V(s_2) = \gamma/(2-\gamma)$.
    \end{enumerate}
  \end{exbox}

  \begin{ideiabox}
    Repare no contraste: o TD usou cada tupla uma vez e propagou pouco;
    o MC propagou o desfecho para todo o episódio; a equivalência de
    certeza extraiu o máximo — ao preço de inverter uma matriz.
  \end{ideiabox}
\end{frame}

\end{document}
